% Formal definitions for Paper L (Manufacturing)
% Included from main.tex

\begin{definition}[アラームトランザクション]
製造装置群から時間ビン幅 $\Delta t$ で収集されたアラームログを
$\mathcal{L} = \{(a_i, t_i)\}_{i=1}^{N}$ とする。
ここで $a_i \in \mathcal{A}$（アラームタイプ集合）、$t_i \in \mathbb{R}^+$（タイムスタンプ）である。
時間ビン $\tau_k = [k\Delta t, (k+1)\Delta t)$ に含まれるアラームタイプの集合を
トランザクション $T_k = \{a_i \mid t_i \in \tau_k\}$ と定義する。
\end{definition}

\begin{definition}[アラームパターンの密集区間]
アラームパターン $P \subseteq \mathcal{A}$ に対し、
ウィンドウサイズ $W$ と閾値 $\theta$ が与えられたとき、
$P$ の\emph{密集区間} $[s, e]$ とは、
任意の $t \in [s, e]$ について
$|\{k \mid P \subseteq T_k, k \in [t, t+W]\}| \geq \theta$
を満たす極大連続区間をいう。
\end{definition}

\begin{definition}[保全イベント帰属]
保全イベント $m$ の実施時点 $t_m$ に対し、
パターン $P$ の密度変化スコアを
$\Delta_P(m) = \bar{s}_P^{\mathrm{post}} - \bar{s}_P^{\mathrm{pre}}$
と定義する。
ここで $\bar{s}_P^{\mathrm{pre}} = \frac{1}{L} \sum_{t=t_m-L}^{t_m-1} s_P(t)$、
$\bar{s}_P^{\mathrm{post}} = \frac{1}{L} \sum_{t=t_m}^{t_m+L-1} s_P(t)$
であり、$s_P(t)$ はパターン $P$ のウィンドウ位置 $t$ におけるサポートカウント、
$L$ はルックバック長である。
\end{definition}

\begin{proposition}[Apriori 性質による候補削減]
アイテムセット $P$ が密集区間を持つならば、
$P$ の任意の部分集合 $P' \subset P$ も密集区間を持つ。
対偶により、密集区間を持たない部分集合が存在する候補は
探索から除外できる。
\end{proposition}
